%%%%%%%%%%%%%%%%%
% This is an sample CV template created using altacv.cls
% (v1.7.2, 28 August 2024) written by LianTze Lim (liantze@gmail.com). Compiles with pdfLaTeX, XeLaTeX and LuaLaTeX.
%
%% It may be distributed and/or modified under the
%% conditions of the LaTeX Project Public License, either version 1.3
%% of this license or (at your option) any later version.
%% The latest version of this license is in
%%    http://www.latex-project.org/lppl.txt
%% and version 1.3 or later is part of all distributions of LaTeX
%% version 2003/12/01 or later.
%%%%%%%%%%%%%%%%

%% Use the "normalphoto" option if you want a normal photo instead of cropped to a circle
% \documentclass[10pt,a4paper,normalphoto]{altacv}

\documentclass[10pt,a4paper,ragged2e,withhyper]{altacv}
%% AltaCV uses the fontawesome5 and simpleicons packages.
%% See http://texdoc.net/pkg/fontawesome5 and http://texdoc.net/pkg/simpleicons for full list of symbols.

% Change the page layout if you need to
\geometry{left=1.2cm,right=1.2cm,top=1.0cm,bottom=1.0cm,columnsep=1.0cm}

% The paracol package lets you typeset columns of text in parallel
\usepackage{paracol}

% Change the font if you want to, depending on whether
% you're using pdflatex or xelatex/lualatex
% WHEN COMPILING WITH XELATEX PLEASE USE
% xelatex -shell-escape -output-driver="xdvipdfmx -z 0" sample.tex
\iftutex 
  % If using xelatex or lualatex:
  \setmainfont{Roboto Slab}
  \setsansfont{Lato}
  \renewcommand{\familydefault}{\sfdefault}
\else
  % If using pdflatex:
  \usepackage[rm]{roboto}
  \usepackage[defaultsans]{lato}
  % \usepackage{sourcesanspro}
  \renewcommand{\familydefault}{\sfdefault}
\fi

% Change the colours if you want to
\definecolor{SlateGrey}{HTML}{2E2E2E}
\definecolor{LightGrey}{HTML}{666666}
\definecolor{DarkPastelRed}{HTML}{450808}
\definecolor{PastelRed}{HTML}{8F0D0D}
\definecolor{GoldenEarth}{HTML}{E7D192}
\colorlet{name}{black}
\colorlet{tagline}{PastelRed}
\colorlet{heading}{DarkPastelRed}
\colorlet{headingrule}{GoldenEarth}
\colorlet{subheading}{PastelRed}
\colorlet{accent}{PastelRed}
\colorlet{emphasis}{SlateGrey}
\colorlet{body}{LightGrey}

% Change some fonts, if necessary
\renewcommand{\namefont}{\Huge\rmfamily\bfseries}
\renewcommand{\personalinfofont}{\footnotesize}
\renewcommand{\cvsectionfont}{\LARGE\rmfamily\bfseries}
\renewcommand{\cvsubsectionfont}{\large\bfseries}


% Change the bullets for itemize and rating marker
% for \cvskill if you want to
\renewcommand{\cvItemMarker}{{\small\textbullet}}
\renewcommand{\cvRatingMarker}{\faCircle}
% ...and the markers for the date/location for \cvevent
% \renewcommand{\cvDateMarker}{\faCalendar*[regular]}
% \renewcommand{\cvLocationMarker}{\faMapMarker*}


% If your CV/résumé is in a language other than English,
% then you probably want to change these so that when you
% copy-paste from the PDF or run pdftotext, the location
% and date marker icons for \cvevent will paste as correct
% translations. For example Spanish:
% \renewcommand{\locationname}{Ubicación}
% \renewcommand{\datename}{Fecha}


%% Use (and optionally edit if necessary) this .tex if you
%% want to use an author-year reference style like APA(6)
%% for your publication list
% % When using APA6 if you need more author names to be listed
% because you're e.g. the 12th author, add apamaxprtauth=12
\usepackage[backend=biber,style=apa6,sorting=ydnt]{biblatex}
\defbibheading{pubtype}{\cvsubsection{#1}}
\renewcommand{\bibsetup}{\vspace*{-\baselineskip}}
\AtEveryBibitem{%
  \makebox[\bibhang][l]{\itemmarker}%
  \iffieldundef{doi}{}{\clearfield{url}}%
}
\setlength{\bibitemsep}{0.25\baselineskip}
\setlength{\bibhang}{1.25em}


%% Use (and optionally edit if necessary) this .tex if you
%% want an originally numerical reference style like IEEE
%% for your publication list
\input{pubs-num.tex}

%% sample.bib contains your publications
\addbibresource{sample.bib}
% \usepackage{academicons}\let\faOrcid\aiOrcid
\begin{document}
\name{Mgr. Norbert Micheľ}
\tagline{Software Developer, Data Analyst}
%% You can add multiple photos on the left or right
%\photoR{2.8cm}{Globe_High}
\photoR{2.8cm}{nieco2}
% \photoL{2.5cm}{Yacht_High,Suitcase_High}

\personalinfo{%
  % Not all of these are required!
  \email{noro.michel159@gmail.com}
  \phone{+421 907 869 883}

  \location{Košice, Slovakia}
  %\homepage{www.homepage.com}
  % \twitter{@twitterhandle}
  %\xtwitter{@x-handle}
  \linkedin{norbert-micheľ}
  \github{nmichel159}
  \printinfo{\faGraduationCap}{Mgr. Charles University}
 
  %\orcid{0000-0000-0000-0000}
  %% You can add your own arbitrary detail with
  %% \printinfo{symbol}{detail}[optional hyperlink prefix]
  % \printinfo{\faPaw}{Hey ho!}[https://example.com/]

  %% Or you can declare your own field with
  %% \NewInfoFiled{fieldname}{symbol}[optional hyperlink prefix] and use it:
  % \NewInfoField{gitlab}{\faGitlab}[https://gitlab.com/]
  % \gitlab{your_id}
  %%
  %% For services and platforms like Mastodon where there isn't a
  %% straightforward relation between the user ID/nickname and the hyperlink,
  %% you can use \printinfo directly e.g.
  % \printinfo{\faMastodon}{@username@instace}[https://instance.url/@username]
  %% But if you absolutely want to create new dedicated info fields for
  %% such platforms, then use \NewInfoField* with a star:
  % \NewInfoField*{mastodon}{\faMastodon}
  %% then you can use \mastodon, with TWO arguments where the 2nd argument is
  %% the full hyperlink.
  % \mastodon{@username@instance}{https://instance.url/@username}
}

\makecvheader
%% Depending on your tastes, you may want to make fonts of itemize environments slightly smaller
% \AtBeginEnvironment{itemize}{\small}

%% Set the left/right column width ratio to 6:4.
\columnratio{0.6}

% Start a 2-column paracol. Both the left and right columns will automatically
% break across pages if things get too long.
\begin{paracol}{2}
\cvsection{Experience}

\cvevent{AI Engineer}{Instaview}{2026 -- Present}{Prague}

\cvevent{Software Developer}{LevelGo}{2025 -- 2026}{Košice}

\cvevent{}{DADO}{2025 -- 2025}{Bratislava}

\cvevent{}{KDC - Košice Development Center}{2021 -- 2024}{Košice}

\divider

\cvevent{Erasmus Student Exchange}{Charles University, Prague}{2024 -- 2024}{Prague}

\divider

\cvevent{Student Scientific Activity}{UPJŠ - P. J. Šafárik University}{2022 -- 2023}{Košice}





\cvsection{Projects}

\cvevent{AI Engineer – Interview Analysis System}{Instaview}{2026}{Prague}
\begin{itemize}
    \item Technologies: \textbf{Python, LLMs}
    \item Developed an AI-driven system for automated interview analysis using large language models.
    \item Implemented semantic evaluation of candidate responses to assess both technical and soft skills.
    \item Designed algorithms for matching candidates to positions based on skill alignment and contextual meaning.
\end{itemize}
\divider

\cvevent{Full-Stack Developer – Mobile Sports Application}{LevelGo}{2025 -- 2026}{Košice}
\begin{itemize}
    \item Technologies: \textbf{Python, Flutter}
    \item Implemented a mobile application using Flutter with a focus on clean architecture and maintainable code structure.
    \item Developed back-end services in Python to support application logic, data processing, and API communication.
\end{itemize}


\divider
\cvevent{Full-Stack Developer – AI Integration Platform}{DADO}{2025 -- 2025}{Bratislava}
\begin{itemize}
    \item Technologies: \textbf{Vue.js, PHP}
    \item Integrated third-party AI APIs for automated data processing and content generation tasks.
    \item Customized and adapted AI tools to meet specific client requirements and workflow automation objectives.
\end{itemize}



\divider


\cvevent{Internal Blazor Component Library}{KDC}{}{}
\begin{itemize}
\item Language: \textbf{C\#}
    \item Designed and implemented a reusable Blazor library that contains custom graphical components such as counters, sliders, and other UI elements.
    \item Developed demo applications showcasing component usage for internal developers.
\end{itemize}
\divider
\cvevent{Full-Stack Application for Slovak Railways - ZSSK}{KDC}{}{}
\begin{itemize}
\item Language: \textbf{C\#}
    \item Developed a full-stack application in C\# for ZSSK, encompassing both front-end and back-end development.
    \item Integrated and synchronized data within the internal ZSSK system.
    \item Managed and resolved bug reports while working on concurrent projects.
\end{itemize}

\divider
\cvevent{Web Front-End Development for Government-Related Websites - Various Clients}{KDC}{}{}
\begin{itemize}
    \item Language - \textbf{JavaScript}
    \item Developed and maintained front-end interfaces for multiple government-related websites, ensuring responsive and user-friendly design.
    \item Implemented and integrated APIs to facilitate seamless communication between front-end and back-end systems for diverse clients.
\end{itemize}

\divider
\cvevent{Application for Scheduling Using Linear Programming}{UPJŠ}{}{}
\begin{itemize}
    \item Language -  \textbf{Python}
    \item Designed and developed an application for scheduling tasks based on linear programming techniques, tailored for the IYPT (International Young Physicists’ Tournament) event.
    \item Addressed the NP-hard nature of the scheduling problem by incorporating elements of randomness to accelerate the algorithm and achieve practical runtimes.
\end{itemize}

\newpage
\cvsection{Achievements}

\cvsubsectionfont{Academic}

\cvevent{Best Scientific Thesis}{Nomination for Bernard Bolzano Foundation Annual Prize, Charles University}{2025}{Prague}
\begin{itemize}
  \item Prestigious annual award recognizing the best scientific thesis at the Faculty of Mathematics and Physics, granted to a very limited number of students.
\end{itemize}

\divider

\cvsubsectionfont{Programming}

\cvevent{Master of Slovakia}{Central European Regional Contest (CERC)}{2022}{Ljubljana}

\cvevent{2nd Place}{Student Programming Competition - UPJŠ}{2022}{Košice}

\cvevent{Winner}{National Round of the Informatics Olympiad (CKOI)}{2020}{Bratislava}

\divider

\cvsubsectionfont{Math}

\cvevent{Silver Medal}{Middle European Mathematical Olympiad (MEMO)}{2020}{Košice}


\medskip

\cvsection{A Day of My Life}

% Adapted from @Jake's answer from http://tex.stackexchange.com/a/82729/226
% \wheelchart{outer radius}{inner radius}{
% comma-separated list of value/text width/color/detail}
\wheelchart{1.5cm}{0.5cm}{%
  8/8em/accent!30/{Sleep},
  1/6em/accent!50/Household tasks,
  8/8em/accent!60/{Daytime job, Academic Life},
  4/10em/accent/Sports and relaxation,
  5/6em/accent!20/Spending time with friends
}

% use ONLY \newpage if you want to force a page break for
% ONLY the current column


%% Switch to the right column. This will now automatically move to the second
%% page if the content is too long.
\switchcolumn
\cvsection{Highlights}



\cvachievement{\faStar}{Master’s Degree in One Year}{Charles University, accelerated completion}
\vspace*{-0.8ex}\divider\vspace*{-0.6ex}
\cvachievement{\faStar}{Silver Medal – MEMO}{International mathematics competition}
\vspace*{-0.8ex}\divider\vspace*{-0.6ex}
\cvachievement{\faStar}{Best Thesis Nomination}{Bernard Bolzano Foundation Prize}
\vspace*{-0.8ex}\divider\vspace*{-0.6ex}
\cvachievement{\faStar}{Slovakia Master – CERC}{Competitive programming champion}

\cvsection{\href{https://github.com/nmichel159}{\faGithub Projects on GitHub}}

\cvevent{\href{https://github.com/nmichel159/3D_tic_tac_toe_dropped}{3D Tic-Tac-Toe Dropped \faGithub}}{Python, Tkinter}{2022}{}
\vspace*{-1ex}\divider\vspace*{-0.8ex}
\cvevent{\href{https://github.com/nmichel159/IYPT_Linear_programming}{IYPT-Linear programming \faGithub}}{Python, Tkinter, PuLP}{2022}{}
\vspace*{-1ex}\divider\vspace*{-0.8ex}
\cvevent{\href{https://github.com/nmichel159/scheduling}{Scheduling \faGithub}}{FastAPI, React, Python}{2026}{}
\vspace*{-1ex}\divider\vspace*{-0.8ex}
\cvevent{\href{https://github.com/nmichel159/skill-analyzer-LLM}{Skill Analyzer LLM \faGithub}}{Python, LLMs, PostgreSQL}{2026}{}
\vspace*{-1ex}\divider\vspace*{-0.8ex}
\cvevent{\href{https://github.com/nmichel159/sport}{Sport App \faGithub}}{FastAPI, React Native, React}{2026}{}
\vspace*{-1ex}\divider\vspace*{-0.8ex}
\cvevent{\href{https://github.com/nmichel159/stock}{Stock Screener \faGithub}}{Python, Streamlit, Pandas}{2026}{}

\cvsection{Languages}

\cvtag{Python}
\cvtag{C++}
\cvtag{C\#}
\cvtag{Haskell}
\cvtag{Java}
\cvtag{JavaScript}
\cvtag{TypeScript}
\cvtag{Rust}
\cvtag{Kotlin}
\cvtag{SQL}

\vspace*{-1ex}
\cvsection{Spoken Languages}
\vspace*{-0.5ex}
\cvskill{English - B2}{4}
\cvskill{Slovak - C1}{5}
\cvskill{Czech - C1}{5}
\cvskill{German - A2}{1}
\newpage

%% Yeah I didn't spend too much time making all the
%% spacing consistent... sorry. Use \smallskip, \medskip,
%% \bigskip, \vspace etc to make adjustments.
\medskip

\cvsection{Education}

\cvevent{PhD. in Computer Science }{Pavol Jozef Šafárik University, Faculty of Science, Košice, Slovakia}{2025 -- present}{}
Thesis title: General Logistic Problems. Supervisor: prof\ RNDr.\ Gabriel Semanišin, PhD.

\divider

\cvevent{Mgr. in Discrete Models and Algorithms}{Charles University, Faculty of Mathematics and Physics, Prague, Czech Republic}{2024 -- 2025}{}
Thesis title: Min Cut-Path Problem. Supervisor: prof.\ RNDr.\ Martin Loebl, CSc.

\divider

\cvevent{Bc. in Data Analysis and Artificial Intelligence}{Pavol Jozef Šafárik University, Faculty of Science, Košice, Slovakia}{2021 -- 2024}{}
Thesis title: Mrežové mnohosteny v n-rozmernej kocke. Supervisor: Mgr.\ Martin Vodička.


\cvsection{Frameworks}

\cvtag{Blazor}
\cvtag{NumPy}
\cvtag{SciPy}
\cvtag{Matplotlib}
\cvtag{Pandas}
\cvtag{NetworkX}
\cvtag{React}
\cvtag{Vue}
\cvtag{Angular}
\cvtag{Flask}
\cvtag{Flutter}

\cvsection{Tools}

\cvtag{LaTeX}
\cvtag{Git}
\cvtag{Jupyter}
\cvtag{Visual Studio}
\cvtag{VS Code}
\cvtag{.NET}


\cvsection{Strengths}

% Don't overuse these \cvtag boxes — they're just eye-candies and not essential. If something doesn't fit on a single line, it probably works better as part of an itemized list (probably inlined itemized list), or just as a comma-separated list of strengths.

% The `ragged2e` document class option might cause automatic linebreaks between \cvtag to fail.
% Either remove the ragged2e option; or 
% add \LaTeXraggedright in the paragraph for these \cvtag
{\LaTeXraggedright
\cvtag{Reliable}
\cvtag{Hard-working}
\cvtag{Flexible}
\cvtag{Friendly}
\cvtag{Adaptable}
\cvtag{Soft-skill}
\par}
\cvsection{Interests}

\cvtag{Abstinent}
\cvtag{Chess}
\cvtag{Ultimate Frisbee}

\newpage

\cvsection{Mentoring}

\cvevent{Construction of a Small-Cap Equity Investment Portfolio Using Large Language Models}{Master's Thesis -- Kiara Koščová}{2026}{}

\divider

\cvevent{Algorithms for Move Generation and Evaluation in Chess}{Bachelor's Thesis -- Tomáš Oltman}{2025}{}

\cvsection{Publications}

%% Specify your last name(s) and first name(s) as given in the .bib to automatically bold your own name in the publications list.
\mynames{Michel/N.,
  Michel/Norbert,
  Micheľ/N.,
  Micheľ/Norbert}

\nocite{*}
\printbibliography[heading=pubtype,title={\printinfo{\faFile*[regular]}{Prepared Publications}},type=unpublished]





% \divider



\end{paracol}


\end{document}
